\begin{lemma}[Elementary sets with at most one colored boundary edge are interchangeable]
Let $G$ be a finite loopless multigraph, let $\varphi$ be a proper partial
$k$-edge-coloring, and let $T$ be a tree-sequence.  Put $X=V(T)$.  If $X$
is elementary under $\varphi$ and, for colors $\alpha,\beta\in[k]$,
\[
 \partial_{\varphi,\alpha}(X)=\varnothing,
 \qquad
 \bigl|\partial_{\varphi,\beta}(X)\bigr|\leq 1,
\]
then $\alpha$ and $\beta$ are $T$-interchangeable under $\varphi$.
The boundary sets and the two-color components retain exact edge identities,
so the assertion also applies in the presence of parallel edges and includes
isolated two-color path components.
\end{lemma}
\begin{proof}
Write
\[
 I=\noderef{TwoColorIncidenceGraph}(G,\varphi;\alpha,\beta),
 \qquad X=\noderef{TreeSequenceVertexSet}(T).
\]
We prove the quantified condition in \noderef{TreeInterchangeable}.
Fix nodes $z,w\in V(G)\sqcup E(G)$ such that their full components
\[
 K_z=\noderef{TwoColorComponent}(G,\varphi,\alpha,\beta,z),
 \qquad
 K_w=\noderef{TwoColorComponent}(G,\varphi,\alpha,\beta,w)
\]
are path components and each contains a vertex-node belonging to $X$.
Unfolding \noderef{TwoColorPathComponent}, either path-component
hypothesis first gives $\alpha\ne\beta$ and supplies in its component a
node whose neighbor set in $I$ has at most one member.

By \noderef{TwoColorComponentsDisjoint}, either $K_z=K_w$, which is the
required conclusion, or $K_z$ and $K_w$ are disjoint.  Suppose for a
contradiction that they are disjoint.  We first classify either one of
these two components; call it $K$.

Apply \noderef{TwoColorComponentSpanningPath} using the assumed
properness.  It gives a simple path $P$ in $I$ whose support is exactly
$K$.  There are two cases.

Suppose first that $K$ contains no edge-node.  Since $K$ meets $X$, it
contains a vertex-node $\mathsf{inl}(x)$ with $x\in X$.  Any adjacency in
$I$ has one vertex-node and one edge-node by
\noderef{TwoColorIncidenceGraph}.  Hence no node of $K$ can be adjacent
to $\mathsf{inl}(x)$.  Since $K$ is a connected component, it follows
that
\[
                         K=\{\mathsf{inl}(x)\}.              \tag{1}
\]
The same adjacency definition says that there is no edge identity
incident with $x$ having color $\alpha$ or $\beta$.  Thus neither color
belongs to \noderef{PresentColors}$(G,\varphi,x)$, and the complementary
definition \noderef{MissingColors} gives
\[
 \alpha,\beta\in\overline\varphi(x).                         \tag{2}
\]
This is exactly the isolated path-component case; it has not been
discarded.

Suppose instead that $K$ contains an edge-node.  Because $K$ also contains
a vertex-node and is connected, every edge-node occurring in $K$ lies on
an incidence walk from a vertex-node.  At the first adjacency to that
edge-node, \noderef{TwoColorIncidenceGraph} says that its edge identity
$f$ has value $\varphi(f)=\alpha$ or $\varphi(f)=\beta$.  Conversely,
both endpoints of such an $f$ are neighbors of its edge-node: incidence
is endpoint membership by \noderef{Incident}.  The two endpoints are
distinct, because equality of them would make the endpoint pair
$G.\operatorname{ends}(f)$ a loop, contrary to
\noderef{LooplessMultigraph}.  Therefore every edge-node in $K$ has
degree exactly two in $I$.                                      \tag{3}

Every node of $I$ has degree at most two by
\noderef{TwoColorIncidenceDegreeAtMostTwo}.  The low-degree node supplied
by \noderef{TwoColorPathComponent} lies on the spanning path $P$ because
the support of $P$ is $K$.  By (3) it is a vertex-node.  It cannot be an
internal node of $P$, since the two distinct neighboring nodes on a simple
path would contradict that its neighbor set is a subsingleton.  Orient
$P$ so this node is its first end.  The other end cannot have an extra
neighbor in $K$: an extra neighbor which is internal on $P$ would give
that internal node three distinct neighbors, while an extra adjacency to
the first end would give the first end two neighbors.  Every neighbor of
a node of $K$ lies again in its full connected component, so there is no
extra neighbor outside $K$.  Consequently both ends of $P$ have degree
one in $I$.  By (3), neither is an edge-node.  Thus for distinct vertices
$r,s$,
\[
       \text{the ends of }P\text{ are }\mathsf{inl}(r),
       \mathsf{inl}(s).                                      \tag{4}
\]
They are distinct because $P$ contains an edge-node and an edge-node has
two distinct endpoint neighbors.

At either end in (4), its unique neighbor is the identity of an incident
edge colored $\alpha$ or $\beta$.  If that identity has color $\alpha$,
then no incident edge can have color $\beta$, since it would give a second
neighbor; hence $\beta$ is absent from \noderef{PresentColors} and belongs
to \noderef{MissingColors}.  The same argument with the colors reversed
applies when the unique edge has color $\beta$.  We have proved that each
of $r,s$ is missing at least one of the two displayed colors.             \tag{5}

We next determine how a nontrivial $K$ can meet $X$.  If all vertex-nodes
of $K$ belong to $X$, call $K$ trapped.  Then (4)--(5) give two distinct
vertices $r,s\in X$, each missing one of $\alpha,\beta$.                  \tag{6}

Otherwise $K$ contains both a vertex of $X$ and a vertex outside $X$.
Traverse $P$ between such vertices.  At the first change from a vertex in
$X$ to a vertex outside $X$, the intervening edge-node is an exact edge
identity $f$ incident with both vertices.  The definition
\noderef{BoundaryEdges}, together with endpoint membership in
\noderef{Incident} and the distinct-end condition of
\noderef{LooplessMultigraph}, gives
\[
                            f\in\partial_G(X).                \tag{7}
\]
The definition of $I$ says $\varphi(f)=\alpha$ or
$\varphi(f)=\beta$.  In the first alternative, (7) would put $f$ in
\noderef{BoundaryColorEdges}$(G,\varphi,X,\alpha)$, contrary to the
assumed emptiness of that finite set.  Therefore
\[
             f\in\partial_{\varphi,\beta}(X).                 \tag{8}
\]
Thus every component which reaches outside $X$ consumes an exact identity
from the $\beta$-colored boundary.  Two disjoint components cannot consume
the same identity, because its edge-node would then lie in both components.
The bound
$|\partial_{\varphi,\beta}(X)|\leq1$ therefore shows that at most one of
$K_z,K_w$ reaches outside $X$.                                \tag{9}

Moreover, a nontrivial component which reaches outside has at least one
of its two ends in $X$.  Indeed, if both ends in (4) were outside while
the path met $X$, traversing the simple path would give a first entrance
into $X$ and a later last departure from $X$.  The two intervening
edge-nodes are distinct because $P$ is simple.  Repeating (7)--(8) for
both would give two distinct identities in
$\partial_{\varphi,\beta}(X)$, contradicting its cardinality bound.
Together with (5), every reaching component therefore supplies a vertex
of $X$ missing $\alpha$ or $\beta$.                           \tag{10}

We now apply the classification to the disjoint $K_z,K_w$.  If one is an
edge-free singleton, let its tree vertex be $x$.  By (2), $x$ misses both
colors.  The other component cannot be another singleton, since its
distinct tree vertex would repeat either missing color.  If it is
nontrivial, it is either trapped and supplies an end in $X$ by (6), or it
reaches outside and supplies such an end by (10).  That end is distinct
from $x$ by disjointness and misses one of $\alpha,\beta$, while $x$
misses the same color.  In every subcase this contradicts elementariness
as defined in \noderef{ColoredVertexSetElementary}.

It remains that both components are nontrivial.  By (9), at least one is
trapped.  Its two distinct ends lie in $X$ and each misses one displayed
color by (6).  The other component supplies at least one end in $X$: use
(6) if it is trapped and (10) if it reaches outside.  Disjointness makes
these three tree vertices distinct.  Label each by one of
$\alpha,\beta$ which it misses, as provided by (5).  Among three vertices
and two labels, two vertices receive the same label.  Their missing-color
sets are not disjoint, contradicting the assumed
\noderef{ColoredVertexSetElementary} property of $X$.

Hence the disjoint alternative is impossible, and $K_z=K_w$.  Since
$z,w$ were arbitrary path components meeting $X$, the full quantified
definition \noderef{TreeInterchangeable} holds.  Throughout the argument
$f$ denotes an element of the original edge type, and component
disjointness is applied to its node $\mathsf{inr}(f)$; parallel edges are
therefore never identified by their common endpoints.
\end{proof}
